\documentclass[10pt]{article}
\usepackage{arydshln}
\usepackage{graphicx}
\usepackage{tcolorbox}
\usepackage[a4paper, margin=1in]{geometry}
\usepackage{amsmath, amsfonts, amssymb}
\usepackage[colorlinks=true,linkcolor=black,citecolor=black,urlcolor=black]{hyperref}
\usepackage{natbib}
\usepackage{url}

% Citing how to
% \nocite{*}
% \bibliographystyle{plain}
% \bibliography{refs}
% \cite{ref of cite}

% ---- Macros ----
\newcommand{\myhrule}[1]{%
  \vspace{#1}%
  \noindent\hrule
  \vspace{#1}%
}
\makeatletter
\newcommand{\bern}{%
  \begin{enumerate}%
  \renewcommand{\labelenumi}{\roman{enumi}.}%
}
\newcommand{\eern}{%
  \end{enumerate}%
}
\makeatother
\newcommand{\FlaThreeByThreeTL}[9]{%
\left(
\begin{array}{cc|c}
#1 & #2 & #3\\
#4 & #5 & #6\\
\hline
#7 & #8 & #9
\end{array}
\right)%
}
\tcbuselibrary{breakable}
\newcommand{\boxedtext}[1]{%
\begin{tcolorbox}[
    colback=white,
    colframe=black,
    sharp corners,
    boxrule=0.5pt,
    breakable
]%
#1
\end{tcolorbox}%
}
\newcommand{\bt}[1]{%
\begin{tcolorbox}[
    colback=white,
    colframe=black,
    sharp corners,
    boxrule=0.5pt,
    breakable
]%
#1
\end{tcolorbox}%
}
\newcommand{\FlaTwoByOne}[2]{%
\left(
\begin{array}{c}
#1\\
\hline
#2
\end{array}
\right)%
}
\newcommand{\FlaThreeByOneB}[3]{%
\left(
\begin{array}{c}
#1\\
\hline
#2\\
#3
\end{array}
\right)%
}
\newcommand{\FlaTwoByTwoSingleLine}[4]{%
\left(
\begin{array}{cc}
#1 & #2\\
#3 & #4
\end{array}
\right)%
}
\newcommand{\complexpm}{\boxed{\pm}}
\newcommand{\FlaOneByTwoSingleLine}[2]{%
\left(
\begin{array}{cc}
    #1 & #2
\end{array}
\right)%
}
\newcommand{\FlaTwoByOneSingleLine}[2]{%
\left(
\begin{array}{cc}
#1\\
#2
\end{array}
\right)%
}
\newcommand{\FlaThreeByThreeBR}[9]{%
\left(
\begin{array}{c|cc}
#1 & #2 & #3 \\ \hline
#4 & #5 & #6 \\
#7 & #8 & #9
\end{array}
\right)%
}
\newcommand{\FlaTwoByTwo}[4]{%
\left(
\begin{array}{c|c}
#1 & #2 \\
\hline
#3 & #4
\end{array}
\right)
}
\newcommand{\FlaOneByThreeR}[3]{%
\left(
\begin{array}{c|cc}
#1 & #2 & #3
\end{array}
\right)%
}
\newcommand{\FlaOneByTwo}[2]{%
\left(
\begin{array}{c|c}
#1 & #2
\end{array}
\right)%
}
\newcommand{\Span}{\operatorname{Span}}
\newcommand{\pbreak}{\vspace{30pt}}
\newcommand{\ud}[1]{\underline{#1}}
\newcommand{\ital}[1]{\textit{#1}}
\newcommand{\fb}[1]{\fbox{#1}}
\newcommand{\vocab}[1]{\fbox{\textbf{\textit{#1}}}}
\newcommand{\eps}{\varepsilon}
\newcommand{\mb}{\mathbf}
\newcommand{\bs}{\boldsymbol}
\newcommand{\mdef}{\overset{\mathrm{def}}{=}}
\newcommand{\h}[1]{\textbf{\fbox{#1}}}
\newcommand{\real}{\mathbb{R}}
\newcommand{\R}{\mathbb{R}}
\newcommand{\C}{\mathbb{C}}
\newcommand{\Null}{\mathcal{N}}
\newcommand{\amp}{&}
\newcommand{\tpose}{^{\top}}
\newcommand{\ident}[1][]{\mathbf{I}_{#1}}
\newcommand{\defeq}{\overset{\cdot}{=}}
\newcommand{\sbsi}[4]{%
	\begin{center}
   		\begin{minipage}{0.5\textwidth}
      			\centering
      			\includegraphics[width=#2\textwidth]{#1}
    		\end{minipage}%
    		\begin{minipage}{0.5\textwidth}
      			\centering
      			\includegraphics[width=#4\textwidth]{#3}
		\end{minipage}
 	\end{center}
}
\newcommand{\ci}[2]{%
  \begin{center}
	\includegraphics[width=#2\textwidth]{#1}
  \end{center}
}
\newcommand{\cii}[4]{%
  \begin{center}
    \includegraphics[width=#2\textwidth]{#1}\hfill
    \includegraphics[width=#4\textwidth]{#3}
  \end{center}
}
\newcommand{\newbp}{%
  \newpage
  \mbox{}
  \newpage
}
\newcommand{\formbox}[1]{%
  \begin{center}
    \fbox{\parbox{0.95\linewidth}{\[
      #1
    \]}}
  \end{center}
}
\newcommand{\euler}{\mathrm{e}}
\newcommand{\expect}{\mathbb{E}}
\newcommand{\bi}{\begin{itemize}}
\newcommand{\ei}{\end{itemize}}
\newcommand{\be}{\begin{enumerate}}
\newcommand{\ee}{\end{enumerate}}
\newcommand{\al}[1]{\ensuremath{\begin{aligned}#1\end{aligned}}}
\DeclareMathOperator*{\argmax}{arg\,max}
\DeclareMathOperator*{\argmin}{arg\,min}
\newcommand{\rightarrowname}[1]{\ensuremath{\xrightarrow{\text{#1}}}}
\newcommand{\leftarrowname}[1]{\ensuremath{\xleftarrow{\text{#1}}}}
\newcommand{\lt}{<}
\newcommand{\gt}{>}
\setlength{\parindent}{0pt}
\newcommand{\tab}{\hspace*{2em}}
\newcommand{\pause}{}
% ------------------------------
\DeclareMathSizes{10}{13}{8}{8}

% optionally:
% \setcounter{secnumdepth}{0} % removes section numbers
\renewcommand{\arraystretch}{1.5}
% \DeclareMathSizes{<text size>}{<math size>}{<script size>}{<scriptscript size>}

\title{PostgreSQL notes}
\author{}
\date{}

\begin{document}

\maketitle

\tableofcontents

\newbp
\section{General}

\subsection{Terminology}

\begin{itemize}
\item \textbf{Table}: A structured collection of data organized into rows and columns. In PostgreSQL, a table represents an entity (e.g., users, orders).

\item \textbf{Row (Tuple)}: A single record in a table. Each row represents one instance of the entity.

\item \textbf{Column (Attribute)}: A named field in a table that represents a specific property of the entity (e.g., age, email). All values in a column share the same data type.

\item \textbf{Schema}: A logical namespace that organizes database objects such as tables, views, and indexes. Helps avoid naming conflicts.

\item \textbf{Database}: A collection of schemas and their objects. It is the top-level container in PostgreSQL.

\item \textbf{Primary Key}: A column (or set of columns) that uniquely identifies each row in a table. Values must be unique and not NULL.

\item \textbf{Foreign Key}: A column (or set of columns) that references a primary key in another table, establishing relationships between tables.

\item \textbf{Constraint}: A rule enforced on a column or table (e.g., NOT NULL, UNIQUE, CHECK) to ensure data integrity.

\item \textbf{Index}: A data structure that improves the speed of data retrieval operations, often at the cost of additional storage and slower writes.

\item \textbf{View}: A virtual table defined by a query. It does not store data itself but presents data from one or more tables.

\item \textbf{Query}: A request for data or information from the database, typically written in SQL.

\item \textbf{NULL}: A special marker representing missing or unknown data (not the same as zero or an empty string).

\item \textbf{Transaction}: A sequence of operations performed as a single unit of work. Transactions follow ACID properties (Atomicity, Consistency, Isolation, Durability).

\item \textbf{Relation}: A formal term for a table in relational databases, based on relational algebra.

\item \textbf{Cardinality}: The number of rows in a table, or the nature of relationships between tables (e.g., one-to-many).

\end{itemize}

\newpage

\subsection{General Structure of a SQL Query}
\begin{itemize}
\item A SQL query follows a logical execution order, which is different from how it is written.
\end{itemize}

\begin{verbatim}
SELECT columns
FROM table
JOIN other_table
ON join_condition
WHERE row_filter
GROUP BY grouping_columns
HAVING group_filter
ORDER BY sorting_columns
LIMIT number_of_rows;
\end{verbatim}

\begin{itemize}
\item Logical execution order:
    \begin{itemize}
    \item FROM (choose tables)
    \item JOIN (combine tables)
    \item WHERE (filter rows)
    \item GROUP BY (group rows)
    \item HAVING (filter groups)
    \item SELECT (choose columns / expressions)
    \item ORDER BY (sort results)
    \item LIMIT (restrict output)
    \end{itemize}
\item Important: The written order is NOT the execution order.
\end{itemize}

\newpage 

\subsection{Query Execution Order}

Although SQL queries are written in a top-down manner, they are logically evaluated in a different order. Understanding this order is important for reasoning about query behavior, especially with filtering and aggregation.

\begin{enumerate}
\item \textbf{FROM}: Determine the source tables and perform joins to construct the initial working set of rows.

\item \textbf{WHERE}: Filter rows based on conditions. This happens before any grouping or aggregation.

\item \textbf{GROUP BY}: Group the remaining rows into sets based on specified columns.

\item \textbf{HAVING}: Filter groups after aggregation (similar to WHERE, but operates on grouped data).

\item \textbf{SELECT}: Compute and return the desired columns or expressions. Aggregations (e.g., COUNT, SUM) are evaluated here.

\item \textbf{DISTINCT}: Remove duplicate rows from the result set (if specified).

\item \textbf{ORDER BY}: Sort the final result set.

\item \textbf{LIMIT / OFFSET}: Restrict the number of rows returned.
\end{enumerate}

\paragraph{Example}
\begin{verbatim}
SELECT department, COUNT(*)
FROM employees
WHERE salary > 50000
GROUP BY department
HAVING COUNT(*) > 5
ORDER BY COUNT(*) DESC
LIMIT 10;
\end{verbatim}

\paragraph{Explanation}
\begin{itemize}
\item First, rows are gathered from \texttt{employees} (FROM)
\item Then filtered by salary (WHERE)
\item Remaining rows are grouped by department (GROUP BY)
\item Groups with too few rows are removed (HAVING)
\item Counts are computed and selected (SELECT)
\item Results are sorted (ORDER BY)
\item Finally, only the top 10 are returned (LIMIT)
\end{itemize}

\newpage

\subsection{TRUNCATE vs DROP vs DELETE}
\begin{itemize}
    \item Purpose: Different ways to remove data or tables in SQL.
\end{itemize}

\textbf{DELETE (remove rows selectively or all rows):}

\begin{verbatim}
DELETE FROM employees;

DELETE FROM employees
WHERE salary < 50000;
\end{verbatim}

\begin{itemize}
    \item Removes rows one by one
    \item Can use WHERE clause
    \item Slower for large tables
    \item Keeps table structure
\end{itemize}

\textbf{TRUNCATE (remove all rows quickly):}

\begin{verbatim}
TRUNCATE TABLE employees;
\end{verbatim}

\begin{itemize}
    \item Removes all rows at once
    \item Cannot use WHERE clause
    \item Faster than DELETE
    \item Keeps table structure
\end{itemize}

\textbf{DROP (remove entire table):}

\begin{verbatim}
DROP TABLE employees;
\end{verbatim}

\begin{itemize}
    \item Removes table and all data permanently
    \item Table no longer exists
    \item Fast operation
    \item Must recreate table if needed again
\end{itemize}

\newpage
\subsection{JOIN Differences}
The matches will fill in a NULL value for entries that don't have a match\\ 
(in the case of like all left or all right).
\begin{itemize}
\item JOIN -- implicitly INNER JOIN (only matches)
\item INNER JOIN -- only matches
\item LEFT JOIN / LEFT OUTER JOIN -- all left + matches
\item RIGHT JOIN / RIGHT OUTER JOIN -- all right + matches
\item FULL JOIN / FULL OUTER JOIN -- all rows from both sides
\item CROSS JOIN -- cartesian product (every row on left matches every row on right)
\item SELF JOIN -- table joined with itself
\item NATURAL JOIN -- auto-join on columns with same name (can be unsafe)
\item JOIN USING -- join using a shared column name
\item LATERAL JOIN -- subquery can reference left table (per-row evaluation)
\item SEMI JOIN (EXISTS) -- rows from left that have a match (no duplication)
\item ANTI JOIN (LEFT + NULL) -- rows from left with no match
\end{itemize}

\subsection{NULL Behavior}
NULL means "unknown value" and behaves differently from normal values.
\begin{itemize}
\item Use IS NULL / IS NOT NULL (NOT =)
\item NULL in comparisons results in UNKNOWN (not true/false)
\item Aggregates ignore NULL values (except COUNT(*))
\end{itemize}

\subsection{WHERE vs HAVING}
\begin{itemize}
\item WHERE filters rows BEFORE grouping
\item HAVING filters AFTER GROUP BY
\end{itemize}

\subsection{GROUP BY Rule}
\begin{itemize}
\item Every selected column must be either:
    \begin{itemize}
    \item in GROUP BY
    \item or an aggregate function (COUNT, MAX, etc.)
    \end{itemize}
\end{itemize}

\subsection{ORDER BY vs GROUP BY}
\begin{itemize}
\item GROUP BY = aggregates rows
\item ORDER BY = sorts final result
\end{itemize}

\subsection{COUNT Differences}
\begin{itemize}
\item COUNT(*) counts all rows
\item COUNT(column) ignores NULL values
\end{itemize}

\subsection{DISTINCT}
\begin{itemize}
\item Removes duplicate rows from result
\item Works on entire selected row combination
\end{itemize}

\subsection{JOIN Condition}
\begin{itemize}
\item ON clause defines how tables match
\item Missing or wrong ON can cause Cartesian product (explosion of rows)
\end{itemize}

\subsection{Cartesian Product}
\begin{itemize}
\item Occurs when JOIN has no condition
\item Every row combines with every other row
\item Avoid unless intentionally using CROSS JOIN
\end{itemize}

\subsection{Aliases}
\begin{itemize}
\item Temporary names for tables or columns
\item Improves readability and reduces repetition
\end{itemize}

\textbf{Table alias:}

\begin{verbatim}
SELECT c.id, c.name
FROM customers c;
\end{verbatim}

\textbf{Table aliases with JOIN:}

\begin{verbatim}
SELECT c.name, o.order_id
FROM customers c
JOIN orders o ON c.id = o.customer_id;
\end{verbatim}

\textbf{Column alias:}

\begin{verbatim}
SELECT salary AS monthly_salary
FROM employees;
\end{verbatim}

\subsection{FROM Table vs Derived Table (Subquery / CTE)}
\begin{itemize}
    \item Purpose: Compare querying a base table directly vs using a temporary derived result.
\end{itemize}

\textbf{Direct table access:}

\begin{verbatim}
SELECT *
FROM employees;
\end{verbatim}

\textbf{Derived table (subquery in FROM clause):}

\begin{verbatim}
SELECT *
FROM (
    SELECT id, salary
    FROM employees
    WHERE salary > 50000
) AS high_paid;
\end{verbatim}

\textbf{CTE equivalent (cleaner version):}

\begin{verbatim}
WITH high_paid AS (
    SELECT id, salary
    FROM employees
    WHERE salary > 50000
)
SELECT *
FROM high_paid;
\end{verbatim}

\subsection{EXISTS and \texttt{SELECT 1}}

In SQL, when using \texttt{EXISTS} or \texttt{NOT EXISTS}, the actual value in the \texttt{SELECT} clause does not matter. The query only checks whether at least one row exists.

Thus, the following are all equivalent:

\begin{verbatim}
SELECT 1
SELECT NULL
SELECT *
SELECT column_name
\end{verbatim}

They all produce the same result because \texttt{EXISTS} ignores the selected values and only evaluates the presence of rows. For this reason, \texttt{SELECT 1} is commonly used for clarity and convention.

\newbp
\section{PostgreSQL Commands (Comprehensive Cheat Sheet)}

% ================= SELECT =================
\subsection{SELECT}
\begin{itemize}
    \item Purpose: Retrieve data from one or more tables.
\end{itemize}

\begin{verbatim}
SELECT * FROM employees;

SELECT name, salary
FROM employees
WHERE salary > 50000;
\end{verbatim}

% ================= WHERE =================
\subsection{WHERE}
\begin{itemize}
    \item Purpose: Filter rows based on conditions.
\end{itemize}

\begin{verbatim}
SELECT *
FROM employees
WHERE department = 'Engineering'
  AND salary > 70000;
\end{verbatim}

% ================= ORDER BY =================
\subsection{ORDER BY}
\begin{itemize}
    \item Purpose: Sort results ascending or descending.
\end{itemize}

\begin{verbatim}
SELECT *
FROM employees
ORDER BY salary DESC;
\end{verbatim}

% ================= DISTINCT =================
\subsection{DISTINCT}
\begin{itemize}
    \item Purpose: Remove duplicates.
\end{itemize}

\begin{verbatim}
SELECT DISTINCT department
FROM employees;
\end{verbatim}

% ================= LIMIT / OFFSET =================
\subsection{LIMIT / OFFSET}
\begin{itemize}
    \item Purpose: Restrict number of rows and paginate results.
\end{itemize}

\begin{verbatim}
SELECT *
FROM employees
LIMIT 10 OFFSET 20;
\end{verbatim}

% ================= GROUP BY =================
\subsection{GROUP BY}
\begin{itemize}
    \item Purpose: Aggregate rows into groups.
\end{itemize}

\begin{verbatim}
SELECT department, COUNT(*)
FROM employees
GROUP BY department;
\end{verbatim}

% ================= HAVING =================
\subsection{HAVING}
\begin{itemize}
    \item Purpose: Filter aggregated groups.
\end{itemize}

\begin{verbatim}
SELECT department, COUNT(*)
FROM employees
GROUP BY department
HAVING COUNT(*) > 5;
\end{verbatim}

% ================= JOINS =================
\subsection{JOIN}
Implicitly an INNER JOIN

\subsection{INNER JOIN}
\begin{itemize}
    \item Purpose: Only matching rows between tables.
\end{itemize}

\begin{verbatim}
SELECT e.name, d.name
FROM employees e
INNER JOIN departments d
ON e.department_id = d.id;
\end{verbatim}

\subsection{LEFT JOIN / LEFT OUTER JOIN}
\begin{itemize}
    \item Purpose: All left table rows + matches from right.
    \item Note: LEFT JOIN and LEFT OUTER JOIN are identical.
\end{itemize}

\begin{verbatim}
SELECT e.name, d.name
FROM employees e
LEFT OUTER JOIN departments d
ON e.department_id = d.id;
\end{verbatim}

\subsection{RIGHT JOIN / RIGHT OUTER JOIN}
\begin{itemize}
    \item Purpose: All right table rows + matches from left.
    \item Note: RIGHT JOIN and RIGHT OUTER JOIN are identical.
\end{itemize}

\begin{verbatim}
SELECT e.name, d.name
FROM employees e
RIGHT OUTER JOIN departments d
ON e.department_id = d.id;
\end{verbatim}

\subsection{FULL OUTER JOIN}
\begin{itemize}
    \item Purpose: All rows from both tables (matched where possible, NULL otherwise).
\end{itemize}

\begin{verbatim}
SELECT e.name, d.name
FROM employees e
FULL OUTER JOIN departments d
ON e.department_id = d.id;
\end{verbatim}

\subsection{CROSS JOIN}
\begin{itemize}
    \item Purpose: Cartesian product (all combinations of rows).
    \item No ON clause is used.
\end{itemize}

\begin{verbatim}
SELECT e.name, s.subject
FROM employees e
CROSS JOIN subjects s;
\end{verbatim}

\subsection{SELF JOIN}
\begin{itemize}
    \item Purpose: Join a table with itself.
    \item Commonly used for hierarchical data (e.g., employee-manager).
\end{itemize}

\begin{verbatim}
SELECT e.name AS employee, m.name AS manager
FROM employees e
LEFT JOIN employees m
ON e.manager_id = m.id;
\end{verbatim}

\subsection{NATURAL JOIN}
\begin{itemize}
    \item Purpose: Automatically joins columns with the same name.
    \item Warning: Can be error-prone if schemas change.
\end{itemize}

\begin{verbatim}
SELECT *
FROM employees
NATURAL JOIN departments;
\end{verbatim}

\subsection{JOIN USING}
\begin{itemize}
    \item Purpose: Simplifies joins when both tables share the same column name.
\end{itemize}

\begin{verbatim}
SELECT *
FROM employees
JOIN departments
USING (department_id);
\end{verbatim}

\subsection{LATERAL JOIN}
\begin{itemize}
    \item Purpose: Allows a subquery to reference columns from the left table.
    \item Very useful for per-row subqueries.
\end{itemize}

\begin{verbatim}
SELECT e.name, sub.*
FROM employees e
LEFT JOIN LATERAL (
    SELECT *
    FROM projects p
    WHERE p.employee_id = e.id
    LIMIT 1
) sub ON true;
\end{verbatim}

\subsection{ANTI JOIN (pattern)}
\begin{itemize}
    \item Purpose: Find rows in one table with no match in another.
    \item Implemented using LEFT JOIN + WHERE NULL.
\end{itemize}

\begin{verbatim}
SELECT e.*
FROM employees e
LEFT JOIN departments d
ON e.department_id = d.id
WHERE d.id IS NULL;
\end{verbatim}

\subsection{SEMI JOIN (pattern)}
\begin{itemize}
    \item Purpose: Return rows that have a match (without duplicating).
    \item Typically written using EXISTS.
\end{itemize}

\begin{verbatim}
SELECT e.*
FROM employees e
WHERE EXISTS (
    SELECT 1
    FROM departments d
    WHERE e.department_id = d.id
);
\end{verbatim}

% ================= INSERT =================
\subsection{INSERT}
\begin{itemize}
    \item Purpose: Insert new rows.
\end{itemize}

\begin{verbatim}
INSERT INTO employees (name, salary)
VALUES ('Alice', 90000);
\end{verbatim}

% ================= UPDATE =================
\subsection{UPDATE}
\begin{itemize}
    \item Purpose: Modify existing rows.
\end{itemize}

\begin{verbatim}
UPDATE employees
SET salary = salary * 1.1
WHERE department = 'Engineering';
\end{verbatim}

% ================= DELETE =================
\subsection{DELETE}
\begin{itemize}
    \item Purpose: Remove rows.
\end{itemize}

\begin{verbatim}
DELETE FROM employees
WHERE id = 10;
\end{verbatim}

% ================= CASE =================
\subsection{CASE WHEN}
\begin{itemize}
    \item Purpose: Conditional logic in queries.
\end{itemize}

\begin{verbatim}
SELECT name,
       CASE
           WHEN salary > 100000 THEN 'High'
           WHEN salary > 50000 THEN 'Medium'
           ELSE 'Low'
       END AS salary_level
FROM employees;
\end{verbatim}

% ================= IN =================
\subsection{IN}
\begin{itemize}
    \item Purpose: Match against multiple values.
\end{itemize}

\begin{verbatim}
SELECT *
FROM employees
WHERE department IN ('Engineering', 'HR');
\end{verbatim}

% ================= BETWEEN =================
\subsection{BETWEEN}
\begin{itemize}
    \item Purpose: Range filtering.
\end{itemize}

\begin{verbatim}
SELECT *
FROM employees
WHERE salary BETWEEN 50000 AND 100000;
\end{verbatim}

% ================= LIKE =================
\subsection{LIKE}
\begin{itemize}
    \item Purpose: Pattern matching.
\end{itemize}

\begin{verbatim}
SELECT *
FROM employees
WHERE name LIKE 'A%';
\end{verbatim}

% ================= NULL CHECK =================
\subsection{NULL CHECK}
\begin{itemize}
    \item Purpose: Check missing values.
\end{itemize}

\begin{verbatim}
SELECT *
FROM employees
WHERE manager_id IS NULL;
\end{verbatim}

% ================= AGGREGATES =================
\subsection{\textbf{Aggregate Functions}}
\begin{itemize}
    \item Purpose: Compute summary statistics.
\end{itemize}

\begin{verbatim}
SELECT
    COUNT(*),
    SUM(salary),
    AVG(salary),
    MAX(salary),
    MIN(salary)
FROM employees;
\end{verbatim}

% ================= WINDOW FUNCTIONS =================
\subsection{Window Functions}
\begin{itemize}
    \item Purpose: Compute values across partitions without collapsing rows.
\end{itemize}

\begin{verbatim}
SELECT name,
       salary,
       RANK() OVER (PARTITION BY department ORDER BY salary DESC)
FROM employees;
\end{verbatim}

% ================= CTE =================
\subsection{WITH (CTE)}
\begin{itemize}
    \item Purpose: Create temporary named result sets.
\end{itemize}

\begin{verbatim}
WITH high_salary AS (
    SELECT *
    FROM employees
    WHERE salary > 100000
)
SELECT * FROM high_salary;
\end{verbatim}

% ================= SUBQUERY =================
\newpage
\subsection{\textbf{Subqueries}}
\begin{itemize}
    \item Purpose: Nest queries inside another query.
\end{itemize}

\begin{verbatim}
SELECT name
FROM employees
WHERE salary > (
    SELECT AVG(salary)
    FROM employees
);
\end{verbatim}

% ================= UNION =================
\subsection{UNION / UNION ALL}
\begin{itemize}
    \item Purpose: Combine results of multiple queries.
\end{itemize}

\begin{verbatim}
SELECT name FROM employees
UNION
SELECT name FROM contractors;
\end{verbatim}

% ================= EXISTS =================
\subsection{EXISTS}
\begin{itemize}
    \item Purpose: Check existence of rows in subquery.
\end{itemize}

\begin{verbatim}
SELECT name
FROM employees e
WHERE EXISTS (
    SELECT 1
    FROM departments d
    WHERE d.id = e.department_id
);
\end{verbatim}

% ================= INDEX =================
\subsection{INDEX}
\begin{itemize}
    \item Purpose: Speed up queries.
\end{itemize}

\begin{verbatim}
CREATE INDEX idx_salary
ON employees(salary);
\end{verbatim}

% ================= TRANSACTIONS =================
\subsection{\textbf{Transactions}}
\begin{itemize}
    \item Purpose: Ensure atomic operations.
\end{itemize}

\begin{verbatim}
BEGIN;

UPDATE employees SET salary = salary * 1.1;

COMMIT;
\end{verbatim}

\newpage 
% ================= ALTER TABLE =================
\subsection{ALTER TABLE}
\begin{itemize}
    \item Purpose: Modify an existing table (add/remove columns, change types).
\end{itemize}

\begin{verbatim}
ALTER TABLE employees
ADD COLUMN age INT;

ALTER TABLE employees
DROP COLUMN age;
\end{verbatim}

% ================= DROP TABLE =================
\subsection{DROP TABLE}
\begin{itemize}
    \item Purpose: Permanently delete a table and all its data.
\end{itemize}

\begin{verbatim}
DROP TABLE employees;
\end{verbatim}

% ================= TRUNCATE =================
\subsection{TRUNCATE}
\begin{itemize}
    \item Purpose: Remove all rows from a table quickly (keeps structure).
\end{itemize}

\begin{verbatim}
TRUNCATE TABLE employees;
\end{verbatim}

% ================= CREATE TABLE =================
\subsection{CREATE TABLE}
\begin{itemize}
    \item Purpose: Create a new table with defined columns.
\end{itemize}

\begin{verbatim}
CREATE TABLE employees (
    id SERIAL PRIMARY KEY,
    name TEXT,
    salary INT
);
\end{verbatim}

% ================= CREATE DATABASE =================
\subsection{CREATE DATABASE}
\begin{itemize}
    \item Purpose: Create a new database.
\end{itemize}

\begin{verbatim}
CREATE DATABASE company_db;
\end{verbatim}

% ================= CREATE VIEW =================
\subsection{CREATE VIEW}
\begin{itemize}
    \item Purpose: Create a virtual table based on a query.
\end{itemize}

\begin{verbatim}
CREATE VIEW high_salary AS
SELECT name, salary
FROM employees
WHERE salary > 100000;
\end{verbatim}

% ================= DROP VIEW =================
\subsection{DROP VIEW}
\begin{itemize}
    \item Purpose: Remove a view.
\end{itemize}

\begin{verbatim}
DROP VIEW high_salary;
\end{verbatim}

% ================= INSERT ... ON CONFLICT =================
\subsection{INSERT INTO ... ON CONFLICT}
\begin{itemize}
    \item Purpose: Handle duplicate key errors during insert.
\end{itemize}

\begin{verbatim}
INSERT INTO users (id, name)
VALUES (1, 'Alice')
ON CONFLICT (id) DO NOTHING;
\end{verbatim}

% ================= UPSERT =================
\subsection{UPSERT (ON CONFLICT DO UPDATE)}
\begin{itemize}
    \item Purpose: Insert a row or update it if a conflict occurs.
\end{itemize}

\begin{verbatim}
INSERT INTO users (id, name)
VALUES (1, 'Alice')
ON CONFLICT (id)
DO UPDATE SET name = EXCLUDED.name;
\end{verbatim}

% ================= MERGE =================
\subsection{MERGE}
\begin{itemize}
    \item Purpose: Conditionally insert, update, or delete based on matches.
\end{itemize}

\begin{verbatim}
MERGE INTO employees e
USING new_data n
ON e.id = n.id
WHEN MATCHED THEN
    UPDATE SET salary = n.salary
WHEN NOT MATCHED THEN
    INSERT (id, salary) VALUES (n.id, n.salary);
\end{verbatim}

% ================= RETURNING =================
\subsection{RETURNING}
\begin{itemize}
    \item Purpose: Return affected rows after INSERT/UPDATE/DELETE.
\end{itemize}

\begin{verbatim}
INSERT INTO users (name)
VALUES ('Alice')
RETURNING id;
\end{verbatim}

% ================= COALESCE =================
\subsection{COALESCE}
\begin{itemize}
    \item Purpose: Return the first non-NULL value.
\end{itemize}

\begin{verbatim}
SELECT COALESCE(phone, 'N/A')
FROM users;
\end{verbatim}

% ================= NULLIF =================
\subsection{NULLIF}
\begin{itemize}
    \item Purpose: Return NULL if two values are equal.
\end{itemize}

\begin{verbatim}
SELECT NULLIF(score, 0)
FROM results;
\end{verbatim}

% ================= CAST =================
\subsection{CAST}
\begin{itemize}
    \item Purpose: Convert a value from one data type to another.
\end{itemize}

\begin{verbatim}
SELECT CAST(salary AS TEXT)
FROM employees;

-- or shorthand
SELECT salary::TEXT
FROM employees;
\end{verbatim}

% ================= :: TYPE CASTING =================
\subsection{:: type casting}
\begin{itemize}
    \item Purpose: Shorthand syntax for casting between data types (PostgreSQL).
\end{itemize}

\begin{verbatim}
SELECT '123'::INT;

SELECT salary::TEXT
FROM employees;
\end{verbatim}

% ================= STRING FUNCTIONS =================
\subsection{STRING FUNCTIONS (CONCAT, SUBSTRING, LENGTH)}
\begin{itemize}
    \item Purpose: Manipulate and analyze text data.
\end{itemize}

\begin{verbatim}
SELECT CONCAT(first_name, ' ', last_name)
FROM employees;

SELECT SUBSTRING(name FROM 1 FOR 3)
FROM employees;

SELECT LENGTH(name)
FROM employees;
\end{verbatim}

% ================= DATE FUNCTIONS =================
\subsection{DATE FUNCTIONS (NOW, CURRENT\_DATE, AGE)}
\begin{itemize}
    \item Purpose: Work with dates and time values.
\end{itemize}

\begin{verbatim}
SELECT NOW();

SELECT CURRENT_DATE;

SELECT AGE(NOW(), birth_date)
FROM users;
\end{verbatim}

% ================= INTERVAL =================
\subsection{INTERVAL}
\begin{itemize}
    \item Purpose: Represent and manipulate durations of time.
\end{itemize}

\begin{verbatim}
SELECT NOW() + INTERVAL '1 day';

SELECT NOW() - INTERVAL '2 hours';

SELECT CURRENT_DATE + INTERVAL '7 days';
\end{verbatim}

% ================= ARRAYS =================
\subsection{ARRAYS}
\begin{itemize}
    \item Purpose: Store and query lists of values in a single column.
\end{itemize}

\begin{verbatim}
CREATE TABLE students (
    id SERIAL,
    grades INT[]
);

INSERT INTO students (grades)
VALUES (ARRAY[90, 85, 88]);

SELECT grades[1]
FROM students;
\end{verbatim}

% ================= UNNEST =================
\subsection{UNNEST}
\begin{itemize}
    \item Purpose: Expand an array into a set of rows.
\end{itemize}

\begin{verbatim}
SELECT UNNEST(ARRAY[1,2,3,4]);

SELECT id, UNNEST(grades)
FROM students;
\end{verbatim}

% ================= JSON / JSONB =================
\subsection{JSON / JSONB OPERATIONS}
\begin{itemize}
    \item Purpose: Store and query structured JSON data.
\end{itemize}

\begin{verbatim}
SELECT data->'name'
FROM users;

SELECT data->>'name'
FROM users;

SELECT data->>'age'
FROM users
WHERE data->>'city' = 'Houston';
\end{verbatim}

% ================= CTE RECURSIVE =================
\subsection{CTE RECURSIVE}
\begin{itemize}
    \item Purpose: Perform recursive queries (e.g., hierarchies, trees).
\end{itemize}

\begin{verbatim}
WITH RECURSIVE nums AS (
    SELECT 1 AS n
    UNION ALL
    SELECT n + 1 FROM nums WHERE n < 5
)
SELECT * FROM nums;
\end{verbatim}

% ================= WINDOW FUNCTION (LEAD) =================
\subsection{WINDOW FUNCTIONS (LEAD)}
\begin{itemize}
    \item Purpose: Access the next row's value without grouping.
\end{itemize}

\begin{verbatim}
SELECT name, salary,
       LEAD(salary) OVER (ORDER BY salary) AS next_salary
FROM employees;
\end{verbatim}

% ================= WINDOW FUNCTION (LAG) =================
\subsection{WINDOW FUNCTIONS (LAG)}
\begin{itemize}
    \item Purpose: Access the previous row's value without grouping.
\end{itemize}

\begin{verbatim}
SELECT name, salary,
       LAG(salary) OVER (ORDER BY salary) AS prev_salary
FROM employees;
\end{verbatim}

% ================= NTILE =================
\subsection{NTILE}
\begin{itemize}
    \item Purpose: Divide rows into equal-sized buckets.
\end{itemize}

\begin{verbatim}
SELECT name, salary,
       NTILE(4) OVER (ORDER BY salary) AS quartile
FROM employees;
\end{verbatim}

% ================= ROW_NUMBER =================
\subsection{ROW\_NUMBER}
\begin{itemize}
    \item Purpose: Assign a unique sequential number to each row.
\end{itemize}

\begin{verbatim}
SELECT name, salary,
       ROW_NUMBER() OVER (ORDER BY salary DESC) AS row_num
FROM employees;
\end{verbatim}

% ================= RANK =================
\subsection{RANK}
\begin{itemize}
    \item Purpose: Rank rows with gaps for ties.
\end{itemize}

\begin{verbatim}
SELECT name, salary,
       RANK() OVER (ORDER BY salary DESC) AS rank
FROM employees;
\end{verbatim}

% ================= DENSE_RANK =================
\subsection{DENSE\_RANK}
\begin{itemize}
    \item Purpose: Rank rows without gaps for ties.
\end{itemize}

\begin{verbatim}
SELECT name, salary,
       DENSE_RANK() OVER (ORDER BY salary DESC) AS dense_rank
FROM employees;
\end{verbatim}

% ================= PARTITION BY =================
\subsection{PARTITION BY}
\begin{itemize}
    \item Purpose: Divide rows into groups for window functions.
\end{itemize}

\begin{verbatim}
SELECT department, name, salary,
       ROW_NUMBER() OVER (PARTITION BY department ORDER BY salary DESC)
FROM employees;
\end{verbatim}

% ================= OVER CLAUSE =================
\subsection{OVER CLAUSE}
\begin{itemize}
    \item Purpose: Define the window for window functions.
\end{itemize}

\begin{verbatim}
SELECT name, salary,
       SUM(salary) OVER () AS total_salary,
       AVG(salary) OVER (PARTITION BY department) AS avg_dept_salary
FROM employees;
\end{verbatim}

% ================= EXPLAIN =================
\subsection{EXPLAIN}
\begin{itemize}
    \item Purpose: Show the query execution plan without running it.
\end{itemize}

\begin{verbatim}
EXPLAIN
SELECT * FROM employees
WHERE salary > 50000;
\end{verbatim}

% ================= EXPLAIN ANALYZE =================
\subsection{EXPLAIN ANALYZE}
\begin{itemize}
    \item Purpose: Execute the query and show the actual execution plan with timing.
\end{itemize}

\begin{verbatim}
EXPLAIN ANALYZE
SELECT * FROM employees
WHERE salary > 50000;
\end{verbatim}

% ================= VACUUM =================
\subsection{VACUUM}
\begin{itemize}
    \item Purpose: Reclaim storage and clean up dead tuples.
\end{itemize}

\begin{verbatim}
VACUUM employees;
\end{verbatim}

% ================= ANALYZE =================
\subsection{ANALYZE}
\begin{itemize}
    \item Purpose: Update table statistics for the query planner.
\end{itemize}

\begin{verbatim}
ANALYZE employees;
\end{verbatim}

% ================= REINDEX =================
\subsection{REINDEX}
\begin{itemize}
    \item Purpose: Rebuild indexes to improve performance or fix corruption.
\end{itemize}

\begin{verbatim}
REINDEX TABLE employees;

REINDEX INDEX idx_salary;
\end{verbatim}

% ================= GRANT =================
\subsection{GRANT}
\begin{itemize}
    \item Purpose: Give permissions to users or roles.
\end{itemize}

\begin{verbatim}
GRANT SELECT, INSERT ON employees TO user1;

GRANT ALL PRIVILEGES ON TABLE employees TO admin_role;
\end{verbatim}

% ================= REVOKE =================
\subsection{REVOKE}
\begin{itemize}
    \item Purpose: Remove permissions from users or roles.
\end{itemize}

\begin{verbatim}
REVOKE SELECT ON employees FROM user1;

REVOKE ALL PRIVILEGES ON TABLE employees FROM admin_role;
\end{verbatim}

% ================= IFNULL =================
\subsection{IFNULL}
\begin{itemize}
    \item Purpose: Returns a replacement value when an expression is \texttt{NULL}.
    \item Use Case: Common in SQL (especially MySQL) to handle missing or undefined values in query results.
    \item Note: In PostgreSQL, \texttt{COALESCE()} is used instead of \texttt{IFNULL()}.
\end{itemize}

\begin{verbatim}
-- MySQL example
SELECT IFNULL(score, 0) AS score
FROM players;

-- Equivalent PostgreSQL version
SELECT COALESCE(score, 0) AS score
FROM players;
\end{verbatim}

% ================= UPPER =================
\subsection{UPPER}
\begin{itemize}
    \item Purpose: Converts all characters in a string to uppercase.
    \item Use Case: Commonly used to make string comparisons or sorting case-insensitive.
\end{itemize}

\begin{verbatim}
SELECT UPPER(name)
FROM users
ORDER BY UPPER(name);
\end{verbatim}

% ================= LOWER =================
\subsection{LOWER}
\begin{itemize}
    \item Purpose: Converts all characters in a string to lowercase.
    \item Use Case: Used for case-insensitive comparisons, filtering, and sorting.
\end{itemize}

\begin{verbatim}
SELECT LOWER(name)
FROM users
ORDER BY LOWER(name);
\end{verbatim}

% ================= COLLATE "C" =================
\subsection{COLLATE "C"}
\begin{itemize}
    \item Purpose: Forces byte-level string comparison instead of locale-aware sorting.
	\item Use Case: Ensures deterministic and system-independent \textbf{lexicographic ordering}.
\end{itemize}

\begin{verbatim}
SELECT name
FROM users
ORDER BY name COLLATE "C";
\end{verbatim}

\begin{itemize}
    \item ``C'' $\rightarrow$ byte-order sorting (fast, deterministic), also described as raw machine-level sorting.
    \item ``POSIX'' $\rightarrow$ essentially equivalent to ``C'' on most systems, providing the same deterministic ordering behavior.
    \item ``en\_US.UTF-8'' $\rightarrow$ applies English language locale rules for human-friendly sorting.
    \item ``fr\_FR.UTF-8'' $\rightarrow$ applies French locale rules, including language-specific character ordering.
\end{itemize}

% ================= STRING_AGG =================
\subsection{STRING\_AGG}
\begin{itemize}
    \item Concatenates values into a single string with a delimiter
    \item Syntax:
\begin{verbatim}
STRING_AGG(column, 'delimiter')
\end{verbatim}
    \item Example:
\begin{verbatim}
SELECT STRING_AGG(product, ', ')
FROM Activities;
\end{verbatim}
    \item With DISTINCT and ordering:
\begin{verbatim}
SELECT STRING_AGG(DISTINCT product, ',' ORDER BY product)
FROM Activities;
\end{verbatim}
    \item Common use: readable output (e.g., comma-separated lists)
\end{itemize}

% ================= ARRAY_AGG =================
\subsection{ARRAY\_AGG}
\begin{itemize}
    \item Aggregates values into a PostgreSQL array
    \item Syntax:
\begin{verbatim}
ARRAY_AGG(column)
\end{verbatim}
    \item Example:
\begin{verbatim}
SELECT ARRAY_AGG(product)
FROM Activities;
\end{verbatim}
    \item With DISTINCT and ordering:
\begin{verbatim}
SELECT ARRAY_AGG(DISTINCT product ORDER BY product)
FROM Activities;
\end{verbatim}
    \item Common use: when you need a true list for further processing
\end{itemize}

% ================= SIMILAR TO =================
\subsection{SIMILAR TO}
\begin{itemize}
    \item Purpose: Pattern matching operator in PostgreSQL that combines aspects of \texttt{LIKE} and regular expressions.
    \item Use Case: Used when simple pattern matching is needed without full regex complexity.
    \item Syntax: Uses SQL-style wildcards (\% and \_) and supports limited regex-like constructs.
\end{itemize}

\begin{verbatim}
SELECT *
FROM Users
WHERE email SIMILAR TO '[A-Za-z][A-Za-z0-9._-]*@leetcode\.com';
\end{verbatim}

\begin{itemize}
    \item Note: SIMILAR TO is stricter than LIKE but less powerful than regex (~).
\end{itemize}

% ================= ~ (REGEX MATCH) =================
\subsection{\textasciitilde{} (Regex Match Operator)}
\begin{itemize}
    \item Purpose: Performs full POSIX regular expression matching in PostgreSQL.
    \item Use Case: Used for advanced string validation, parsing, and pattern enforcement.
    \item Operators:
    \begin{itemize}
        \item \texttt{\textasciitilde{}} → matches regex (case-sensitive)
        \item \texttt{\textasciitilde*} → matches regex (case-insensitive)
        \item \texttt{!\textasciitilde} → does NOT match regex
        \item \texttt{!\textasciitilde{}*} → does NOT match regex (case-insensitive)
    \end{itemize}
\end{itemize}

\begin{verbatim}
SELECT user_id, name, mail
FROM Users
WHERE mail ~ '^[A-Za-z][A-Za-z0-9_.-]*@leetcode\.com$';
\end{verbatim}

\begin{itemize}
    \item Note: Preferred over SIMILAR TO when complex validation rules are required.
\end{itemize}

More examples: 
\begin{itemize}
    \item \texttt{\textasciitilde{}} → matches regex (case-sensitive)

\begin{verbatim}
WHERE mail ~ '^[A-Za-z][A-Za-z0-9_.-]*@leetcode\.com$'
\end{verbatim}

    \item \texttt{\textasciitilde*} → matches regex (case-insensitive)

\begin{verbatim}
WHERE mail ~* '^[a-z][a-z0-9_.-]*@leetcode\.com$'
\end{verbatim}

    \item \texttt{!\textasciitilde} → does NOT match regex

\begin{verbatim}
WHERE mail !~ '^[A-Za-z][A-Za-z0-9_.-]*@leetcode\.com$'
\end{verbatim}

    \item \texttt{!\textasciitilde{}*} → does NOT match regex (case-insensitive)

\begin{verbatim}
WHERE mail !~* '^[a-z][a-z0-9_.-]*@leetcode\.com$'
\end{verbatim}
\end{itemize}

\end{document}


